%%=============================================================================
%% Samenvatting
%%=============================================================================

% TODO: De "abstract" of samenvatting is een kernachtige (~ 1 blz. voor een
% thesis) synthese van het document.
%
% Een goede abstract biedt een kernachtig antwoord op volgende vragen:
%
% 1. Waarover gaat de bachelorproef?
% 2. Waarom heb je er over geschreven?
% 3. Hoe heb je het onderzoek uitgevoerd?
% 4. Wat waren de resultaten? Wat blijkt uit je onderzoek?
% 5. Wat betekenen je resultaten? Wat is de relevantie voor het werkveld?
%
% Daarom bestaat een abstract uit volgende componenten:
%
% - inleiding + kaderen thema
% - probleemstelling
% - (centrale) onderzoeksvraag
% - onderzoeksdoelstelling
% - methodologie
% - resultaten (beperk tot de belangrijkste, relevant voor de onderzoeksvraag)
% - conclusies, aanbevelingen, beperkingen
%
% LET OP! Een samenvatting is GEEN voorwoord!

%%---------- Nederlandse samenvatting -----------------------------------------
%
% TODO: Als je je bachelorproef in het Engels schrijft, moet je eerst een
% Nederlandse samenvatting invoegen. Haal daarvoor onderstaande code uit
% commentaar.
% Wie zijn bachelorproef in het Nederlands schrijft, kan dit negeren, de inhoud
% wordt niet in het document ingevoegd.

\IfLanguageName{english}{%
\selectlanguage{dutch}
\chapter*{Samenvatting}
\lipsum[1-4]
\selectlanguage{english}
}{}

%%---------- Samenvatting -----------------------------------------------------
% De samenvatting in de hoofdtaal van het document

\chapter*{\IfLanguageName{dutch}{Samenvatting}{Abstract}}

Digitale gebruikers beschikken over voldoende digitale basisvaardigheden, maar onderzoek toont aan dat zij moeilijk te overtuigen zijn om wachtwoordmanagers te gebruiken. Ondanks hun vertrouwdheid met computers ervaren zij drempels zoals wantrouwen, beperkte kennis van voordelen en onzekerheid over veiligheid en gebruiksgemak. Aangezien veilig wachtwoordbeheer noodzakelijk is om digitale risico’s te beperken, bestaat er nood aan een leeromgeving waarin gebruikers veilig kunnen oefenen met het gebruik van een wachtwoordmanager.
Deze bachelorproef onderzoekt in welke mate een game-engine, in het bijzonder Unity, geschikt is voor het ontwikkelen van een webgebaseerde simulatie waarin beginnende gebruikers op een toegankelijke en veilige manier kunnen kennismaken met de functionaliteit van wachtwoordmanagers. Het onderzoek vertrekt vanuit bestaande literatuur over de obstakels die zij ondervinden bij het gebruik van een wachtwoordmanager en welke risico's ze lopen bij het niet gebruiken hiervan. Op basis van deze inzichten worden de behoeften en gebruiksproblemen van de doelgroep geanalyseerd, samen met de technische en didactische vereisten voor een veilige oefenomgeving. 
Vervolgens wordt een proof-of-concept uitgewerkt, die kan worden ingezet als workshop, waarin gebruikers in een veilige sandbox omgeving kunnen oefenen met het gebruik van een wachtwoordmanager. In deze simulatie worden realistische scenario's gesimuleerd zonder gebruik te maken van echte persoonlijke gegevens. 
De geschiktheid van Unity wordt beoordeeld op vlak van toegankelijkheid, veiligheid, performance en bruikbaarheid in webomgevingen. De evaluatie richt zich op de ontwikkelbaarheid van dergelijke simulaties, de meerwaarde voor de doelgroep en de beperkingen van Unity, in vergelijking met alternatieven, zoals bestaande kennisclips of webtoepassingen. De resultaten bieden inzicht in de haalbaarheid en de pedagogische waarde van webgebaseerde simulaties voor het versterken van online veilig gedrag bij beginnende gebruikers zonder ervaring met wachtwoordmanagers.

\section{Het gebruik van AI in de bachelorproef}

In deze bachelorproef werd ChatGPT \url{https://www.chatgpt.com} gebruikt om tot een beter begrip van inzichten te komen in de vakliteratuur, alsook bij het oplossen van problemen met Latex (invoegen van afbeeldingen, compilatieproblemen).
Bij het maken van de technische component, de POC, werd de app Claude gebruikt \url{https://claude.com/download}. Deze AI app heeft mij geholpen bij het schrijven van C\+ code in Unity. 
